\chapter{18-09-26}

\section{Máximo común divisor: definición y propiedades}
Para $a,b\in\mathbb Z$, consideramos
\[
\operatorname{Div}_{\geq0}(a,b)=\{c\in\mathbb N_0:c\mid a,\ c\mid b\}.
\]
Si alguno de $a,b$ es no nulo, este conjunto es no vacío y finito: contiene a $1$ y sus elementos están acotados por el valor absoluto de cualquier componente no nula.
\begin{definicion}[{title={Máximo común divisor}}]
Para $a,b$ no simultáneamente nulos, definimos
\[
\operatorname{mcd}(a,b)=\max\operatorname{Div}_{\geq0}(a,b).
\]
Además, adoptamos $\operatorname{mcd}(0,0)=0$. También usamos $(a:b)$ para el máximo común divisor.
\end{definicion}
\begin{proposicion}
Para $a,b\in\mathbb Z$:
\begin{enumerate}
\item $\operatorname{mcd}(a,b)=\operatorname{mcd}(b,a)$.
\item $\operatorname{mcd}(a,0)=|a|$.
\item $\operatorname{mcd}(a,b)=\operatorname{mcd}(-a,b)=\operatorname{mcd}(a,-b)=\operatorname{mcd}(|a|,|b|)$.
\item $\operatorname{mcd}(a,b)=|a|$ si y solo si $a\mid b$.
\end{enumerate}
\end{proposicion}
\begin{demostracion}
La simetría y la invariancia por signos se siguen de la igualdad de los conjuntos de divisores comunes. Si $a\neq0$, los divisores comunes no negativos de $a$ y $0$ son los divisores positivos de $a$, cuyo máximo es $|a|$. Para $a=0$ vale la convención adoptada.
Si $a\mid b$ y $a\neq0$, $|a|$ es divisor común y ningún divisor positivo de $a$ lo supera. Recíprocamente, si el máximo común divisor es $|a|$, entonces $|a|\mid b$ y, por tanto, $a\mid b$. Para $a=0$, ambos enunciados equivalen a $b=0$.
\end{demostracion}

\subsection{Invariancia al restar un múltiplo}
\begin{proposicion}
Para $a,b,c\in\mathbb Z$,
\[
\operatorname{mcd}(a,b)=\operatorname{mcd}(a,b-ca).
\]
\end{proposicion}
\begin{demostracion}
Para cualquier entero $t$,
\[
t\mid a,\ t\mid b\quad\Longleftrightarrow\quad t\mid a,\ t\mid(b-ca).
\]
La implicación directa usa la resta; la recíproca usa $b=(b-ca)+ca$. Luego los conjuntos de divisores comunes no negativos coinciden. Si $(a,b)\neq(0,0)$, también $(a,b-ca)\neq(0,0)$: de $a=0$ y $b-ca=0$ se deduce $b=0$. Los máximos existen y son iguales. Si $a=b=0$, se aplica la convención $\operatorname{mcd}(0,0)=0$.
\end{demostracion}
\begin{teoria}[{title={Algoritmo de Euclides}}]
Si $a>0$ y $b=ca+r$, con $c\in\mathbb Z$ y $0\leq r<a$, entonces
\[
\operatorname{mcd}(a,b)=\operatorname{mcd}(a,r).
\]
Se repite la división, intercambiando los argumentos cuando sea necesario. Los restos no nulos disminuyen estrictamente; el último resto no nulo es el máximo común divisor.
\end{teoria}
\begin{ejemplo}
Como $102=4\cdot24+6$ y $24=4\cdot6$, resulta
\[
\operatorname{mcd}(102,24)=\operatorname{mcd}(6,24)=\operatorname{mcd}(6,0)=6.
\]
\end{ejemplo}

\section{Combinaciones lineales e identidad de Bézout}
\begin{proposicion}
Sean $a,b\in\mathbb Z$ no simultáneamente nulos. El conjunto
\[
S(a,b)=\{xa+yb:x,y\in\mathbb Z,\ xa+yb>0\}
\]
es un subconjunto no vacío de $\mathbb N$ y
\[
\min S(a,b)=\operatorname{mcd}(a,b).
\]
\end{proposicion}
\begin{demostracion}
Todo elemento de $S(a,b)$ es un entero positivo. Si $a\neq0$, $|a|\in S(a,b)$: se toma $x=1$ o $x=-1$ según su signo, e $y=0$. Si $a=0$, entonces $b\neq0$ y, de igual manera, $|b|\in S(a,b)$. Por tanto, $S(a,b)$ no es vacío.

Por el principio del buen orden existe $d=\min S(a,b)>0$. Existen $x,y\in\mathbb Z$ con $d=xa+yb$. Veamos que $d\mid a$. Dividimos $a$ por $d$:
\[
a=qd+r,\qquad 0\leq r<d.
\]
Si $r>0$, entonces
\[
r=a-qd=a-q(xa+yb)=(1-qx)a+(-qy)b\in S(a,b),
\]
lo cual contradice $r<d=\min S(a,b)$. Luego $r=0$ y $d\mid a$. Dividiendo $b$ por $d$ y repitiendo el argumento, se obtiene $d\mid b$. Así, $d$ es divisor común positivo y $d\leq\operatorname{mcd}(a,b)$.

Sea $e=\operatorname{mcd}(a,b)>0$. Como $e\mid a$ y $e\mid b$, se tiene $e\mid xa+yb=d$. Puesto que $d>0$, resulta $e\leq d$. Ambas desigualdades dan $d=e$.
\end{demostracion}
\begin{proposicion}[{title={Identidad de Bézout}}]
Para todos $a,b\in\mathbb Z$ existen $x,y\in\mathbb Z$ tales que
\[
xa+yb=\operatorname{mcd}(a,b).
\]
\end{proposicion}
\begin{demostracion}
Si $(a,b)\neq(0,0)$, el mínimo anterior pertenece a $S(a,b)$. Si $a=b=0$, basta tomar $x=y=0$.
\end{demostracion}
\begin{ejemplo}
Para $a=3$ y $b=2$:
\[
1\cdot3+(-1)\cdot2=1,\qquad (-1)\cdot3+2\cdot2=1.
\]
Los coeficientes de Bézout no son necesariamente únicos.
\end{ejemplo}
\begin{ejemplo}
\textbf{Retroceder en el algoritmo de Euclides.} Para $11$ y $3$:
\[
11=3\cdot3+2,\qquad3=1\cdot2+1,\qquad2=2\cdot1+0.
\]
Luego $\operatorname{mcd}(11,3)=1$ y
\[
1=3-2=3-(11-3\cdot3)=4\cdot3-11.
\]
Así, $x=-1$ e $y=4$ satisfacen $11x+3y=1$.
\end{ejemplo}

\section{Caracterización y corolarios del máximo común divisor}
\begin{proposicion}
Sean $a,b\in\mathbb Z$. El número $d=\operatorname{mcd}(a,b)$ es el único entero no negativo que satisface:
\begin{enumerate}
\item $d\mid a$ y $d\mid b$;
\item todo divisor común de $a$ y $b$ divide a $d$.
\end{enumerate}
\end{proposicion}
\begin{demostracion}
El máximo común divisor divide a ambos números, incluso para $(a,b)=(0,0)$, pues $0\mid0$. Si $e\mid a$ y $e\mid b$, Bézout da $e\mid xa+yb=d$.
Para la unicidad, si $d'$ también cumple las propiedades, entonces $d'\mid d$ y $d\mid d'$. Dos enteros no negativos que se dividen mutuamente son iguales; luego $d=d'$.
\end{demostracion}
\begin{proposicion}[{title={Corolario: compatibilidad con la divisibilidad}}]
Si $a\mid a'$ y $b\mid b'$, entonces
\[
\operatorname{mcd}(a,b)\mid\operatorname{mcd}(a',b').
\]
\end{proposicion}
\begin{demostracion}
Si $e=\operatorname{mcd}(a,b)$, entonces $e\mid a\mid a'$ y $e\mid b\mid b'$. Por tanto $e$ es divisor común de $a',b'$ y divide a su máximo común divisor.
\end{demostracion}
\begin{proposicion}[{title={Corolario: multiplicación por un factor común}}]
Para $a,b,c\in\mathbb Z$,
\[
\operatorname{mcd}(ac,bc)=|c|\operatorname{mcd}(a,b).
\]
\end{proposicion}
\begin{demostracion}
Primero supongamos $c>0$ y $(a,b)\neq(0,0)$. Sean $d=\operatorname{mcd}(a,b)$ y $e=\operatorname{mcd}(ac,bc)$. Por Bézout existen $x,y\in\mathbb Z$ con $d=xa+yb$; luego
\[
dc=x(ac)+y(bc),
\]
y por tanto $e\mid dc$. Por otro lado, $d\mid a$ y $d\mid b$ implican $dc\mid ac$ y $dc\mid bc$, así que $dc\mid e$. Como $e,dc>0$, se concluye $e=dc$.
Si $c<0$, la invariancia por signos permite reemplazarlo por $|c|$. Si $c=0$ o $a=b=0$, ambos miembros son cero.
\end{demostracion}
\begin{nota}{Precisión sobre el signo}
En el pizarrón aparece la fórmula con el factor $c$. Esa versión requiere $c\geq0$. Para admitir cualquier entero $c$, el factor correcto es $|c|$.
\end{nota}

\section{Coprimalidad y números primos}
\begin{definicion}[{title={Enteros coprimos}}]
Dos enteros $a,b$ son \textbf{coprimos} si $\operatorname{mcd}(a,b)=1$.
\end{definicion}
\begin{nota}{Bézout y coprimalidad}
Como consecuencia de Bézout, $a,b$ son coprimos si y solo si existen enteros $x,y$ con $xa+yb=1$. Para la recíproca, cualquier divisor común divide a $1$, por lo que el máximo común divisor es $1$.
\end{nota}
\begin{definicion}[{title={Número primo}}]
Un entero $p\geq2$ es \textbf{primo} si sus únicos divisores positivos son $1$ y $p$. Un entero $n\geq2$ que no es primo es \textbf{compuesto}.
\end{definicion}
\begin{proposicion}[{title={Lema}}]
Sea $p$ primo y $a\in\mathbb Z$. Entonces
\[
\operatorname{mcd}(a,p)=\begin{cases}1,&p\nmid a,\\p,&p\mid a.\end{cases}
\]
\end{proposicion}
\begin{demostracion}
El máximo común divisor es un divisor positivo de $p$: es $1$ o $p$. Es $p$ exactamente cuando $p$ divide también a $a$.
\end{demostracion}
\begin{proposicion}[{title={Lema: existencia de un divisor primo}}]
Todo entero $n\geq2$ es primo o es divisible por un primo menor que él. En particular, todo entero $n\geq2$ tiene un divisor primo.
\end{proposicion}
\begin{demostracion}
Hacemos inducción fuerte. Para $n=2$ vale porque $2$ es primo. Sea $m>2$ y supongamos el resultado para $2,\dots,m-1$.
Si $m$ es primo, ya está. Si no, tiene un divisor $k$ con $1<k<m$. Por la hipótesis inductiva, o bien $k$ es primo, o bien hay un primo $p<k$ que lo divide. En el primer caso, $k$ es un divisor primo de $m$ menor que $m$. En el segundo, $p\mid k\mid m$ y $p<m$. Esto prueba el resultado.
\end{demostracion}
\subsection{La criba de Eratóstenes}
Se listan los enteros desde $2$, se conserva el menor aún no tachado y se tachan sus múltiplos mayores que él. Se repite con el siguiente número no tachado. En el ejemplo del pizarrón, entre $1$ y $24$, los primos son
\[
2,3,5,7,11,13,17,19,23.
\]
El $1$ no es primo. Todo compuesto se elimina porque tiene un divisor primo menor que él; ningún primo se elimina como múltiplo propio de otro entero mayor que $1$.

\section{Teorema de Euclides: infinitud de los primos}
\begin{proposicion}[{title={Teorema de Euclides}}]
Existen infinitos números primos.
\end{proposicion}
\begin{demostracion}
Supongamos que hay solamente una cantidad finita. Sean $p_1<\cdots<p_M$ todos ellos y consideremos
\[
A=p_1p_2\cdots p_M+1\geq2.
\]
Por el lema anterior, $A$ es primo o tiene un divisor primo.
Si $A$ es primo, debería coincidir con alguno de los $p_i$, pero $A>p_M$, lo que es imposible. Si $A$ es compuesto, tiene un divisor primo $p_i$ de la lista. Entonces $p_i\mid A$ y $p_i\mid p_1\cdots p_M$, por lo que
\[
p_i\mid A-p_1\cdots p_M=1.
\]
Esto contradice $p_i\geq2$. En ambos casos hay contradicción; por tanto, existen infinitos primos.
\end{demostracion}
\begin{nota}{El producto más uno no siempre es primo}
La prueba no afirma que $p_1\cdots p_M+1$ sea siempre primo. Si es compuesto, alguno de sus divisores primos queda fuera de la lista propuesta; eso basta para la contradicción.
\end{nota}

\section{Otros ejercicios fotografiados ese día}
\subsection{Ejercicio de parcial: valores posibles de un mcd}
\begin{ejercicio}
Sean $a,b\in\mathbb Z$ tales que $(a:b)=7$ y $b\equiv1\pmod3$. Hallar los posibles valores de $\operatorname{mcd}(63a-b^2,42)$ y dar un ejemplo para cada caso.
\end{ejercicio}
\begin{solucion}
La resolución siguiente completa el enunciado fotografiado. Como $7\mid a$ y $7\mid b$, se tiene $7\mid63a-b^2$. Además,
\[
63a-b^2\equiv-1\equiv2\pmod3,
\]
pues $b\equiv1\pmod3$. Por tanto $3\nmid63a-b^2$.
Los divisores positivos de $42=2\cdot3\cdot7$ que son múltiplos de $7$ y no de $3$ son $7$ y $14$. Ambos se alcanzan:
\[
\begin{array}{c|c|c|c}
a&b&63a-b^2&\operatorname{mcd}(63a-b^2,42)\\\hline
14&7&833&7\\7&7&392&14
\end{array}
\]
En ambos casos $(a:b)=7$ y $b\equiv1\pmod3$. Los valores posibles son $\boxed{7\text{ y }14}$.
\end{solucion}

\subsection{Fragmento de un ejercicio de divisibilidad}
En dos fotos aparece el siguiente desarrollo, sin que se vea el enunciado completo. Para $a,b\in\mathbb Z\setminus\{0\}$:
\begin{align*}
\frac{3(a+b)}{a}+\frac6b\in\mathbb Z
&\Longleftrightarrow\frac{3(a+b)b+6a}{ab}\in\mathbb Z\\
&\Longleftrightarrow\frac{3b^2+6a}{ab}\in\mathbb Z\\
&\Longleftrightarrow ab\mid3b^2+6a.
\end{align*}
En el segundo paso se restó el entero $3$. De la última condición se deduce $a\mid3b^2$ y $b\mid6a$.
\begin{nota}{Hipótesis que no se puede omitir}
En general, exigir que $a$ y $b$ dividan a un número no garantiza que su producto lo divida. La equivalencia
\[
ab\mid N\quad\Longleftrightarrow\quad a\mid N\text{ y }b\mid N
\]
es válida si $a,b$ son coprimos. Sin esa hipótesis, la implicación recíproca puede fallar: $2\mid2$ y $2\mid2$, pero $4\nmid2$. Como no se ve el enunciado, no se presupone que $a,b$ sean coprimos ni se inventa una clasificación de soluciones.
\end{nota}
